\documentclass[UTF8,fontset=none,11pt,a4paper]{ctexart}
\usepackage{ipara-handbook}
\handbooksetup{DS-B03 Auxiliary Structure × Graph Algorithm}{408 · 数据结构 Bridge B03}{Operation Handoff · Invariant Progress}
\begin{document}
\handbooktitle{辅助结构与图算法}{算法目标怎样调用辅助结构推进不变量}{408 · 数据结构 Internal Bridge B03}

\begin{corebox}{桥梁接口}
\[
\text{Algorithm Needs Operation}\to\text{Data Structure Provides Operation}\to\text{Candidate State Changes}\to\text{Algorithm Invariant Progresses}.
\]
Heap 与 Union-Find 不因为被图算法调用就失去自身 Owner；图算法只获得优先队列的极值操作或并查集的连通判定。
\end{corebox}
\begin{methodbox}{职责分离}
DS05 Owns 上滤、下滤和堆序；DS06 Owns Find、Union 和路径压缩；DS08 Owns cut/relax/入度等全局正确性。调用方只依赖稳定接口，不复制底层实现。
\end{methodbox}
\tableofcontents\newpage

\section{三个交接例子}
Prim 需要从跨越当前树边界的候选边中取最小值；Dijkstra 需要从暂定距离中取最小值并允许新候选入队；Kruskal 需要判断边两端是否已在同一分量，并在安全边加入后合并。三个算法都可使用辅助结构，但“何时安全”不是由 Heap/Union-Find 决定。

\begin{center}\small
\begin{tabularx}{\linewidth}{L{2.4cm}L{3.0cm}YY}\toprule
算法 & 调用接口 & 算法不变量 & 失败/边界\\\midrule
Prim & push/top/pop & 已选顶点集的 cut 安全边 & 非连通图\\
Dijkstra & push/top/pop & 非负权下取出距离已确定 & 负边\\
Kruskal & connected/unite & 已选边无环且逐步扩展森林 & 重复边/非连通\\\bottomrule\end{tabularx}\end{center}

\section{最小调用契约}
\begin{lstlisting}[language=C++]
while (!pending.empty()) {
  Candidate c = pending.pop_min();
  if (algorithm_invariant_allows(c)) commit(c);
  else discard(c);
  add_new_candidates(c);
}
\end{lstlisting}
Kruskal 的对应动作是 `if (!connected(u,v)) { unite(u,v); commit(edge); }`。这里的 `commit` 仍由 DS08 证明；本桥只规定调用方向和信息边界。

\section{做题控制协议}
先圈算法全局目标，再圈辅助结构实际提供的一个操作；将“辅助结构内部复杂度”和“算法候选状态数量”分开相乘。若解答把完整堆代码或并查集递归塞进图算法证明，删回接口级说明并补全算法不变量。
\begin{corebox}{压缩复原}
三问：算法要什么操作？辅助结构保证什么？这个操作怎样让全局不变量前进一步？
\end{corebox}
\end{document}
